% Sekcja architektury — dołącz w main.tex: % Sekcja architektury — dołącz w main.tex: % Sekcja architektury — dołącz w main.tex: % Sekcja architektury — dołącz w main.tex: \input{architektura}
% Wymaga pakietów z main.tex: tcolorbox, graphicx, booktabs, tabularx, enumitem

\section*{Architektura systemów RAG}

W ramach benchmarku porównano trzy podejścia do odpowiadania na pytania dotyczące rekrutacji i studiów na AGH.
Wszystkie warianty korzystają z tej samej bazy wiedzy (dokumenty PDF w Chroma oraz graf Neo4j), lecz różnią się
sposobem doboru kontekstu, liczbą wywołań modelu językowego oraz mechanizmami kontroli jakości odpowiedzi.

\begin{tcolorbox}[colback=blue!5!white,colframe=blue!75!black,title=Zakres porównania]
    \begin{itemize}[leftmargin=*, nosep]
        \item \textbf{Standardowy RAG} --- wyszukiwanie wyłącznie wektorowe (Chroma) + jedno wywołanie LLM.
        \item \textbf{GraphRAG (Hybrid)} --- równoległy retrieval z grafu wiedzy i wektorów + jedno wywołanie LLM.
        \item \textbf{RAG agentowy} --- orkiestracja wieloetapowa (LangGraph) z warstwą bezpieczeństwa, filtrem kontekstu i weryfikacją faktów.
    \end{itemize}
\end{tcolorbox}

\vspace{1em}

% ---------------------------------------------------------------------------
\subsection*{Standardowy RAG (wektorowy)}

\begin{tcolorbox}[colback=gray!5!white,colframe=gray!60!black,title=Przepływ danych --- pipeline liniowy]
    \centering
    \texttt{Zapytanie} $\rightarrow$ \texttt{Embedding} $\rightarrow$ \texttt{Chroma (top-$k$)} $\rightarrow$ \texttt{LLM} $\rightarrow$ \texttt{Odpowiedź}
\end{tcolorbox}

\noindent
Standardowy RAG implementuje klasyczny schemat \textit{Retrieve--Augment--Generate}.
Zapytanie użytkownika jest embedowane modelem \texttt{nomic-embed-text} i porównywane z fragmentami
PDF-ów zgromadzonymi w bazie Chroma. Trzy najbardziej podobne semantycznie chunki trafiają do jednego
promptu, na podstawie którego model generuje odpowiedź.

\begin{itemize}[leftmargin=*]
    \item \textbf{Zalety:} prostota, niska latencja (jedno wywołanie LLM), dobre wyniki przy pytaniach,
          gdy właściwy fragment tekstu leży blisko zapytania w przestrzeni embeddingów.
    \item \textbf{Ograniczenia:} brak rozumienia struktury danych (tabele, progi, relacje kierunek$\rightarrow$przedmiot),
          podatność na \textit{semantic mismatch} --- trafienie w dokument o podobnym słownictwie, lecz innej treści
          (np.\ punkty ECTS zamiast progów rekrutacyjnych).
    \item \textbf{Brak kontroli:} system nie weryfikuje, czy pobrany kontekst jest istotny, ani czy odpowiedź
          jest pokryta źródłami --- model może halucynować lub uczciwie zgłosić brak danych, lecz bez gwarancji spójności.
\end{itemize}

\vspace{1em}

% ---------------------------------------------------------------------------
\subsection*{GraphRAG (Hybrid)}

\begin{tcolorbox}[colback=gray!5!white,colframe=gray!60!black,title=Przepływ danych --- pipeline liniowy z dwoma źródłami]
    \centering
    \texttt{Zapytanie} $\rightarrow$ \texttt{Analiza intencji} $\rightarrow$
    \begin{tabular}{c}
        $\swarrow$ \texttt{Graf Neo4j} \\
        $\searrow$ \texttt{Chroma (wektory)}
    \end{tabular}
    $\rightarrow$ \texttt{LLM} $\rightarrow$ \texttt{Odpowiedź}
\end{tcolorbox}

\noindent
Pipeline hybrydowy rozszerza standardowy RAG o warstwę strukturalną.
Najpierw moduł \textit{intent analyzer} klasyfikuje zapytanie (kategoria: terminy, progi, przedmioty maturalne itd.)
i normalizuje encje (nazwy kierunków, tryby studiów). Na tej podstawie uruchamiane są dwa niezależne wyszukiwania:
zapytania Cypher do grafu Neo4j (dane tabelaryczne, relacje między encjami) oraz wyszukiwanie wektorowe po PDF-ach.

\begin{itemize}[leftmargin=*]
    \item \textbf{Zalety:} precyzyjne zapytania strukturalne (filtrowanie numeryczne, mapowanie kierunek$\rightarrow$atrybut),
          łączenie danych z wielu tabel i dokumentów w jednej odpowiedzi, lepsze wyniki przy pytaniach relacyjnych
          (np.\ dyplom zawodowy $\rightarrow$ wskaźnik P2 na danym kierunku).
    \item \textbf{Ograniczenia:} nadal jedno wywołanie generujące --- model otrzymuje \textit{surowy} kontekst z obu źródeł
          bez selekcji; błędne lub nadmiarowe fragmenty mogą zafałszować odpowiedź.
          Graf może gubić logikę tekstową (np.\ operator \textit{albo} w tabeli przedmiotów maturalnych).
    \item \textbf{Brak kontroli:} podobnie jak w RAG wektorowym, nie ma warstwy bezpieczeństwa ani weryfikacji
          pokrycia odpowiedzi w kontekście.
\end{itemize}

\vspace{1em}

% ---------------------------------------------------------------------------
\subsection*{RAG agentowy (LangGraph)}

\begin{tcolorbox}[colback=purple!5!white,colframe=purple!60!black,title=Przepływ danych --- graf stanów z rozgałęzieniami]
    \centering
    \texttt{Zapytanie} $\rightarrow$ \texttt{Guard} $\rightarrow$ \texttt{Intent} $\rightarrow$ \texttt{Hybrid retrieve}
    $\rightarrow$ \texttt{Context filter} $\rightarrow$ \texttt{Answer} $\rightarrow$ \texttt{Fact check}
    $\rightarrow$ \texttt{Odpowiedź / Clarify / Refuse}
\end{tcolorbox}

\noindent
System agentowy nie jest pojedynczym łańcuchem RAG, lecz \textbf{grafem stanów} zbudowanym w LangGraph.
Każdy węzeł to wyspecjalizowany agent (osobne wywołanie LLM lub logika routingu), a krawędzie warunkowe
decydują o dalszym przebiegu przetwarzania. Dzięki temu pipeline może przerwać generowanie, odmówić
odpowiedzi lub poprosić o doprecyzowanie --- zamiast zawsze zwracać tekst wygenerowany na ślepo.

\begin{figure}[H]
    \centering
    \includegraphics[width=0.72\textwidth]{images/agentic-pipeline.png}
    \caption{Diagram aktywności pipeline'u agentowego. Źródło: implementacja w \texttt{agents/graph.py}.}
\end{figure}

\begin{tcolorbox}[colback=blue!5!white,colframe=blue!75!black,title=Opis etapów pipeline'u agentowego]
    \begin{enumerate}[leftmargin=*, nosep]
        \item \textbf{Guarding Agent (strażnik bezpieczeństwa).}
              Pierwszy punkt kontroli. Agent ocenia, czy zapytanie dotyczy spraw uczelnianych (rekrutacja, studia, AGH),
              czy stanowi próbę \textit{jailbreaku}, \textit{prompt injection} albo temat całkowicie niezwiązany z uczelnią.
              \begin{itemize}[nosep]
                  \item \textit{Bezpieczne} $\rightarrow$ przepływ kontynuowany.
                  \item \textit{Jailbreak / injection / off-topic} $\rightarrow$ węzeł \textbf{Refuse} i zakończenie z komunikatem odmowy.
              \end{itemize}

        \item \textbf{Intent Analyzer (analiza intencji).}
              Ten sam moduł normalizacji co w Hybrid RAG: wyciąga kategorię pytania (terminy, progi, matura, kierunek itd.)
              i standaryzuje nazwy kierunków oraz trybów studiów. Wynik kieruje zapytania do grafu Neo4j.

        \item \textbf{Hybrid-retrieve (pobranie hybrydowe).}
              Równoległe pobranie węzłów z grafu wiedzy oraz fragmentów PDF z Chroma (domyślnie do 5 chunków wektorowych
              plus wyniki grafowe). Na tym etapie \textit{nie} następuje jeszcze generacja odpowiedzi --- zbierany jest
              surowy zestaw kandydatów-kontekstów.

        \item \textbf{Context Filtering Agent (filtr kontekstu).}
              LLM ocenia \textbf{każdy} pobrany fragment pod kątem istotności względem pytania i rozpoznanej intencji.
              Fragmenty oznaczone jako \texttt{reject} są odrzucane (np.\ cykle rekrutacyjne 2 i 3, gdy pytanie dotyczy cyklu 1).
              \begin{itemize}[nosep]
                  \item Brak fragmentów po filtracji $\rightarrow$ \textbf{Clarify template} (prośba o doprecyzowanie / informacja o braku danych).
                  \item Pozostały istotne fragmenty $\rightarrow$ przejście do generacji.
              \end{itemize}

        \item \textbf{Answer Generating Agent (generacja odpowiedzi).}
              Model tworzy \textit{szkic odpowiedzi} (\texttt{draft\_answer}) wyłącznie na podstawie przefiltrowanego kontekstu.
              To nadal etap generatywny, lecz oparty na znacznie ciaśniejszym, jakościowo ocenionym zbiorze źródeł.

        \item \textbf{Fact Check Agent (weryfikacja faktów).}
              Osobny agent sprawdza, czy \textbf{każde istotne twierdzenie} w szkicu ma pokrycie w dostarczonym kontekście.
              Nie ocenia stylu ani kompletności bazy --- jedynie wykrywa halucynacje i niespójności liczbowe.
              \begin{itemize}[nosep]
                  \item \textit{Grounded = true} $\rightarrow$ szkic staje się odpowiedzią końcową (\textbf{Send answer}).
                  \item \textit{Grounded = false} $\rightarrow$ powrót do \textbf{Clarify template} z opisem wykrytych problemów
                        (np.\ błędna suma punktów sportowych).
              \end{itemize}
    \end{enumerate}
\end{tcolorbox}

\vspace{0.5em}

\noindent
Kluczowa różnica względem pozostałych architektur: system agentowy stosuje \textbf{wielokrotne, celowe wywołania LLM}
(zwykle 4--6 na jedno pytanie) z jawnymi punktami decyzyjnymi. Zamiast jednego monolitycznego promptu otrzymujemy
łańcuch wyspecjalizowanych ról: strażnik, selektor kontekstu, generator, audytor. Latencja rośnie, lecz rośnie też
możliwość odrzucenia błędnej odpowiedzi zanim trafi do użytkownika.

\vspace{1em}

% ---------------------------------------------------------------------------
\subsection*{Porównanie architektur}

\noindent
\begin{tabularx}{\textwidth}{@{} >{\raggedright\arraybackslash}p{2.4cm} >{\raggedright\arraybackslash}X >{\raggedright\arraybackslash}X >{\raggedright\arraybackslash}X @{}}
    \toprule
    \textbf{Aspekt} & \textbf{Standardowy RAG} & \textbf{GraphRAG (Hybrid)} & \textbf{RAG agentowy} \\
    \midrule
    Źródła danych &
    Tylko PDF (Chroma) &
  Graf Neo4j + PDF &
  Graf Neo4j + PDF \\
    \addlinespace
    Liczba wywołań LLM &
    1 &
  1 &
  4--6 (guard, filter, answer, fact-check) \\
    \addlinespace
    Analiza intencji &
    Nie &
  Tak (reguły + normalizacja) &
  Tak (ten sam moduł) \\
    \addlinespace
    Filtrowanie kontekstu &
    Brak (top-$k$ embeddingów) &
  Brak (cały retrieval trafia do promptu) &
  Tak (LLM ocenia każdy fragment) \\
    \addlinespace
    Weryfikacja odpowiedzi &
    Brak &
  Brak &
  Tak (fact-check przed wysłaniem) \\
    \addlinespace
    Bezpieczeństwo &
    Brak &
  Brak &
  Tak (guard: jailbreak, injection, off-topic) \\
    \addlinespace
    Ścieżki wyjścia &
    Zawsze odpowiedź LLM &
  Zawsze odpowiedź LLM &
  Odpowiedź / odmowa / prośba o doprecyzowanie \\
    \addlinespace
    Latencja &
    Najniższa &
  Średnia &
  Najwyższa \\
    \bottomrule
\end{tabularx}

\vspace{1em}

\begin{tcolorbox}[colback=green!5!white,colframe=green!50!black,title=Wnioski projektowe]
    \begin{itemize}[leftmargin=*]
        \item \textbf{Komplementarność retrieval i kontroli:} Hybrid RAG rozwiązuje problem \textit{dostępu} do właściwych danych
              (struktura grafu + semantyka wektorów), natomiast RAG agentowy dodaje warstwę \textit{jakości i bezpieczeństwa} ---
              selekcję kontekstu, audyt odpowiedzi oraz blokadę zapytań spoza domeny.
        \item \textbf{Koszt za jakość:} Każdy dodatkowy agent to osobne wywołanie modelu. W benchmarku latencja pipeline'u agentowego
              jest wielokrotnie wyższa niż w wariantach liniowych, co jest akceptowalnym kompromisem w scenariuszach,
              gdzie priorytetem jest wiarygodność (rekrutacja, progi punktowe, regulaminy).
        \item \textbf{Fail-safe zamiast fail-open:} Gdy filtr nie znajdzie istotnych fragmentów lub fact-check wykryje halucynację,
              system \textbf{nie zgaduje} --- zwraca szablon doprecyzowania lub komunikat o braku pokrycia w bazie wiedzy.
              To zachowanie widoczne m.in.\ przy pytaniach o punkty sportowe, gdzie agent odrzucił błędny szkic zamiast
              podać niepoprawną maksymalną liczbę punktów.
        \item \textbf{Ochrona przed atakami:} Warstwa Guard skutecznie blokuje zapytania typu jailbreak i prompt injection
              (kategoria \texttt{security} w zestawie testowym), czego liniowe pipeline'y w ogóle nie implementują.
    \end{itemize}
\end{tcolorbox}

\vspace{1em}

\newpage

% Wymaga pakietów z main.tex: tcolorbox, graphicx, booktabs, tabularx, enumitem

\section*{Architektura systemów RAG}

W ramach benchmarku porównano trzy podejścia do odpowiadania na pytania dotyczące rekrutacji i studiów na AGH.
Wszystkie warianty korzystają z tej samej bazy wiedzy (dokumenty PDF w Chroma oraz graf Neo4j), lecz różnią się
sposobem doboru kontekstu, liczbą wywołań modelu językowego oraz mechanizmami kontroli jakości odpowiedzi.

\begin{tcolorbox}[colback=blue!5!white,colframe=blue!75!black,title=Zakres porównania]
    \begin{itemize}[leftmargin=*, nosep]
        \item \textbf{Standardowy RAG} --- wyszukiwanie wyłącznie wektorowe (Chroma) + jedno wywołanie LLM.
        \item \textbf{GraphRAG (Hybrid)} --- równoległy retrieval z grafu wiedzy i wektorów + jedno wywołanie LLM.
        \item \textbf{RAG agentowy} --- orkiestracja wieloetapowa (LangGraph) z warstwą bezpieczeństwa, filtrem kontekstu i weryfikacją faktów.
    \end{itemize}
\end{tcolorbox}

\vspace{1em}

% ---------------------------------------------------------------------------
\subsection*{Standardowy RAG (wektorowy)}

\begin{tcolorbox}[colback=gray!5!white,colframe=gray!60!black,title=Przepływ danych --- pipeline liniowy]
    \centering
    \texttt{Zapytanie} $\rightarrow$ \texttt{Embedding} $\rightarrow$ \texttt{Chroma (top-$k$)} $\rightarrow$ \texttt{LLM} $\rightarrow$ \texttt{Odpowiedź}
\end{tcolorbox}

\noindent
Standardowy RAG implementuje klasyczny schemat \textit{Retrieve--Augment--Generate}.
Zapytanie użytkownika jest embedowane modelem \texttt{nomic-embed-text} i porównywane z fragmentami
PDF-ów zgromadzonymi w bazie Chroma. Trzy najbardziej podobne semantycznie chunki trafiają do jednego
promptu, na podstawie którego model generuje odpowiedź.

\begin{itemize}[leftmargin=*]
    \item \textbf{Zalety:} prostota, niska latencja (jedno wywołanie LLM), dobre wyniki przy pytaniach,
          gdy właściwy fragment tekstu leży blisko zapytania w przestrzeni embeddingów.
    \item \textbf{Ograniczenia:} brak rozumienia struktury danych (tabele, progi, relacje kierunek$\rightarrow$przedmiot),
          podatność na \textit{semantic mismatch} --- trafienie w dokument o podobnym słownictwie, lecz innej treści
          (np.\ punkty ECTS zamiast progów rekrutacyjnych).
    \item \textbf{Brak kontroli:} system nie weryfikuje, czy pobrany kontekst jest istotny, ani czy odpowiedź
          jest pokryta źródłami --- model może halucynować lub uczciwie zgłosić brak danych, lecz bez gwarancji spójności.
\end{itemize}

\vspace{1em}

% ---------------------------------------------------------------------------
\subsection*{GraphRAG (Hybrid)}

\begin{tcolorbox}[colback=gray!5!white,colframe=gray!60!black,title=Przepływ danych --- pipeline liniowy z dwoma źródłami]
    \centering
    \texttt{Zapytanie} $\rightarrow$ \texttt{Analiza intencji} $\rightarrow$
    \begin{tabular}{c}
        $\swarrow$ \texttt{Graf Neo4j} \\
        $\searrow$ \texttt{Chroma (wektory)}
    \end{tabular}
    $\rightarrow$ \texttt{LLM} $\rightarrow$ \texttt{Odpowiedź}
\end{tcolorbox}

\noindent
Pipeline hybrydowy rozszerza standardowy RAG o warstwę strukturalną.
Najpierw moduł \textit{intent analyzer} klasyfikuje zapytanie (kategoria: terminy, progi, przedmioty maturalne itd.)
i normalizuje encje (nazwy kierunków, tryby studiów). Na tej podstawie uruchamiane są dwa niezależne wyszukiwania:
zapytania Cypher do grafu Neo4j (dane tabelaryczne, relacje między encjami) oraz wyszukiwanie wektorowe po PDF-ach.

\begin{itemize}[leftmargin=*]
    \item \textbf{Zalety:} precyzyjne zapytania strukturalne (filtrowanie numeryczne, mapowanie kierunek$\rightarrow$atrybut),
          łączenie danych z wielu tabel i dokumentów w jednej odpowiedzi, lepsze wyniki przy pytaniach relacyjnych
          (np.\ dyplom zawodowy $\rightarrow$ wskaźnik P2 na danym kierunku).
    \item \textbf{Ograniczenia:} nadal jedno wywołanie generujące --- model otrzymuje \textit{surowy} kontekst z obu źródeł
          bez selekcji; błędne lub nadmiarowe fragmenty mogą zafałszować odpowiedź.
          Graf może gubić logikę tekstową (np.\ operator \textit{albo} w tabeli przedmiotów maturalnych).
    \item \textbf{Brak kontroli:} podobnie jak w RAG wektorowym, nie ma warstwy bezpieczeństwa ani weryfikacji
          pokrycia odpowiedzi w kontekście.
\end{itemize}

\vspace{1em}

% ---------------------------------------------------------------------------
\subsection*{RAG agentowy (LangGraph)}

\begin{tcolorbox}[colback=purple!5!white,colframe=purple!60!black,title=Przepływ danych --- graf stanów z rozgałęzieniami]
    \centering
    \texttt{Zapytanie} $\rightarrow$ \texttt{Guard} $\rightarrow$ \texttt{Intent} $\rightarrow$ \texttt{Hybrid retrieve}
    $\rightarrow$ \texttt{Context filter} $\rightarrow$ \texttt{Answer} $\rightarrow$ \texttt{Fact check}
    $\rightarrow$ \texttt{Odpowiedź / Clarify / Refuse}
\end{tcolorbox}

\noindent
System agentowy nie jest pojedynczym łańcuchem RAG, lecz \textbf{grafem stanów} zbudowanym w LangGraph.
Każdy węzeł to wyspecjalizowany agent (osobne wywołanie LLM lub logika routingu), a krawędzie warunkowe
decydują o dalszym przebiegu przetwarzania. Dzięki temu pipeline może przerwać generowanie, odmówić
odpowiedzi lub poprosić o doprecyzowanie --- zamiast zawsze zwracać tekst wygenerowany na ślepo.

\begin{figure}[H]
    \centering
    \includegraphics[width=0.72\textwidth]{images/agentic-pipeline.png}
    \caption{Diagram aktywności pipeline'u agentowego. Źródło: implementacja w \texttt{agents/graph.py}.}
\end{figure}

\begin{tcolorbox}[colback=blue!5!white,colframe=blue!75!black,title=Opis etapów pipeline'u agentowego]
    \begin{enumerate}[leftmargin=*, nosep]
        \item \textbf{Guarding Agent (strażnik bezpieczeństwa).}
              Pierwszy punkt kontroli. Agent ocenia, czy zapytanie dotyczy spraw uczelnianych (rekrutacja, studia, AGH),
              czy stanowi próbę \textit{jailbreaku}, \textit{prompt injection} albo temat całkowicie niezwiązany z uczelnią.
              \begin{itemize}[nosep]
                  \item \textit{Bezpieczne} $\rightarrow$ przepływ kontynuowany.
                  \item \textit{Jailbreak / injection / off-topic} $\rightarrow$ węzeł \textbf{Refuse} i zakończenie z komunikatem odmowy.
              \end{itemize}

        \item \textbf{Intent Analyzer (analiza intencji).}
              Ten sam moduł normalizacji co w Hybrid RAG: wyciąga kategorię pytania (terminy, progi, matura, kierunek itd.)
              i standaryzuje nazwy kierunków oraz trybów studiów. Wynik kieruje zapytania do grafu Neo4j.

        \item \textbf{Hybrid-retrieve (pobranie hybrydowe).}
              Równoległe pobranie węzłów z grafu wiedzy oraz fragmentów PDF z Chroma (domyślnie do 5 chunków wektorowych
              plus wyniki grafowe). Na tym etapie \textit{nie} następuje jeszcze generacja odpowiedzi --- zbierany jest
              surowy zestaw kandydatów-kontekstów.

        \item \textbf{Context Filtering Agent (filtr kontekstu).}
              LLM ocenia \textbf{każdy} pobrany fragment pod kątem istotności względem pytania i rozpoznanej intencji.
              Fragmenty oznaczone jako \texttt{reject} są odrzucane (np.\ cykle rekrutacyjne 2 i 3, gdy pytanie dotyczy cyklu 1).
              \begin{itemize}[nosep]
                  \item Brak fragmentów po filtracji $\rightarrow$ \textbf{Clarify template} (prośba o doprecyzowanie / informacja o braku danych).
                  \item Pozostały istotne fragmenty $\rightarrow$ przejście do generacji.
              \end{itemize}

        \item \textbf{Answer Generating Agent (generacja odpowiedzi).}
              Model tworzy \textit{szkic odpowiedzi} (\texttt{draft\_answer}) wyłącznie na podstawie przefiltrowanego kontekstu.
              To nadal etap generatywny, lecz oparty na znacznie ciaśniejszym, jakościowo ocenionym zbiorze źródeł.

        \item \textbf{Fact Check Agent (weryfikacja faktów).}
              Osobny agent sprawdza, czy \textbf{każde istotne twierdzenie} w szkicu ma pokrycie w dostarczonym kontekście.
              Nie ocenia stylu ani kompletności bazy --- jedynie wykrywa halucynacje i niespójności liczbowe.
              \begin{itemize}[nosep]
                  \item \textit{Grounded = true} $\rightarrow$ szkic staje się odpowiedzią końcową (\textbf{Send answer}).
                  \item \textit{Grounded = false} $\rightarrow$ powrót do \textbf{Clarify template} z opisem wykrytych problemów
                        (np.\ błędna suma punktów sportowych).
              \end{itemize}
    \end{enumerate}
\end{tcolorbox}

\vspace{0.5em}

\noindent
Kluczowa różnica względem pozostałych architektur: system agentowy stosuje \textbf{wielokrotne, celowe wywołania LLM}
(zwykle 4--6 na jedno pytanie) z jawnymi punktami decyzyjnymi. Zamiast jednego monolitycznego promptu otrzymujemy
łańcuch wyspecjalizowanych ról: strażnik, selektor kontekstu, generator, audytor. Latencja rośnie, lecz rośnie też
możliwość odrzucenia błędnej odpowiedzi zanim trafi do użytkownika.

\vspace{1em}

% ---------------------------------------------------------------------------
\subsection*{Porównanie architektur}

\noindent
\begin{tabularx}{\textwidth}{@{} >{\raggedright\arraybackslash}p{2.4cm} >{\raggedright\arraybackslash}X >{\raggedright\arraybackslash}X >{\raggedright\arraybackslash}X @{}}
    \toprule
    \textbf{Aspekt} & \textbf{Standardowy RAG} & \textbf{GraphRAG (Hybrid)} & \textbf{RAG agentowy} \\
    \midrule
    Źródła danych &
    Tylko PDF (Chroma) &
  Graf Neo4j + PDF &
  Graf Neo4j + PDF \\
    \addlinespace
    Liczba wywołań LLM &
    1 &
  1 &
  4--6 (guard, filter, answer, fact-check) \\
    \addlinespace
    Analiza intencji &
    Nie &
  Tak (reguły + normalizacja) &
  Tak (ten sam moduł) \\
    \addlinespace
    Filtrowanie kontekstu &
    Brak (top-$k$ embeddingów) &
  Brak (cały retrieval trafia do promptu) &
  Tak (LLM ocenia każdy fragment) \\
    \addlinespace
    Weryfikacja odpowiedzi &
    Brak &
  Brak &
  Tak (fact-check przed wysłaniem) \\
    \addlinespace
    Bezpieczeństwo &
    Brak &
  Brak &
  Tak (guard: jailbreak, injection, off-topic) \\
    \addlinespace
    Ścieżki wyjścia &
    Zawsze odpowiedź LLM &
  Zawsze odpowiedź LLM &
  Odpowiedź / odmowa / prośba o doprecyzowanie \\
    \addlinespace
    Latencja &
    Najniższa &
  Średnia &
  Najwyższa \\
    \bottomrule
\end{tabularx}

\vspace{1em}

\begin{tcolorbox}[colback=green!5!white,colframe=green!50!black,title=Wnioski projektowe]
    \begin{itemize}[leftmargin=*]
        \item \textbf{Komplementarność retrieval i kontroli:} Hybrid RAG rozwiązuje problem \textit{dostępu} do właściwych danych
              (struktura grafu + semantyka wektorów), natomiast RAG agentowy dodaje warstwę \textit{jakości i bezpieczeństwa} ---
              selekcję kontekstu, audyt odpowiedzi oraz blokadę zapytań spoza domeny.
        \item \textbf{Koszt za jakość:} Każdy dodatkowy agent to osobne wywołanie modelu. W benchmarku latencja pipeline'u agentowego
              jest wielokrotnie wyższa niż w wariantach liniowych, co jest akceptowalnym kompromisem w scenariuszach,
              gdzie priorytetem jest wiarygodność (rekrutacja, progi punktowe, regulaminy).
        \item \textbf{Fail-safe zamiast fail-open:} Gdy filtr nie znajdzie istotnych fragmentów lub fact-check wykryje halucynację,
              system \textbf{nie zgaduje} --- zwraca szablon doprecyzowania lub komunikat o braku pokrycia w bazie wiedzy.
              To zachowanie widoczne m.in.\ przy pytaniach o punkty sportowe, gdzie agent odrzucił błędny szkic zamiast
              podać niepoprawną maksymalną liczbę punktów.
        \item \textbf{Ochrona przed atakami:} Warstwa Guard skutecznie blokuje zapytania typu jailbreak i prompt injection
              (kategoria \texttt{security} w zestawie testowym), czego liniowe pipeline'y w ogóle nie implementują.
    \end{itemize}
\end{tcolorbox}

\vspace{1em}

\newpage

% Wymaga pakietów z main.tex: tcolorbox, graphicx, booktabs, tabularx, enumitem

\section*{Architektura systemów RAG}

W ramach benchmarku porównano trzy podejścia do odpowiadania na pytania dotyczące rekrutacji i studiów na AGH.
Wszystkie warianty korzystają z tej samej bazy wiedzy (dokumenty PDF w Chroma oraz graf Neo4j), lecz różnią się
sposobem doboru kontekstu, liczbą wywołań modelu językowego oraz mechanizmami kontroli jakości odpowiedzi.

\begin{tcolorbox}[colback=blue!5!white,colframe=blue!75!black,title=Zakres porównania]
    \begin{itemize}[leftmargin=*, nosep]
        \item \textbf{Standardowy RAG} --- wyszukiwanie wyłącznie wektorowe (Chroma) + jedno wywołanie LLM.
        \item \textbf{GraphRAG (Hybrid)} --- równoległy retrieval z grafu wiedzy i wektorów + jedno wywołanie LLM.
        \item \textbf{RAG agentowy} --- orkiestracja wieloetapowa (LangGraph) z warstwą bezpieczeństwa, filtrem kontekstu i weryfikacją faktów.
    \end{itemize}
\end{tcolorbox}

\vspace{1em}

% ---------------------------------------------------------------------------
\subsection*{Standardowy RAG (wektorowy)}

\begin{tcolorbox}[colback=gray!5!white,colframe=gray!60!black,title=Przepływ danych --- pipeline liniowy]
    \centering
    \texttt{Zapytanie} $\rightarrow$ \texttt{Embedding} $\rightarrow$ \texttt{Chroma (top-$k$)} $\rightarrow$ \texttt{LLM} $\rightarrow$ \texttt{Odpowiedź}
\end{tcolorbox}

\noindent
Standardowy RAG implementuje klasyczny schemat \textit{Retrieve--Augment--Generate}.
Zapytanie użytkownika jest embedowane modelem \texttt{nomic-embed-text} i porównywane z fragmentami
PDF-ów zgromadzonymi w bazie Chroma. Trzy najbardziej podobne semantycznie chunki trafiają do jednego
promptu, na podstawie którego model generuje odpowiedź.

\begin{itemize}[leftmargin=*]
    \item \textbf{Zalety:} prostota, niska latencja (jedno wywołanie LLM), dobre wyniki przy pytaniach,
          gdy właściwy fragment tekstu leży blisko zapytania w przestrzeni embeddingów.
    \item \textbf{Ograniczenia:} brak rozumienia struktury danych (tabele, progi, relacje kierunek$\rightarrow$przedmiot),
          podatność na \textit{semantic mismatch} --- trafienie w dokument o podobnym słownictwie, lecz innej treści
          (np.\ punkty ECTS zamiast progów rekrutacyjnych).
    \item \textbf{Brak kontroli:} system nie weryfikuje, czy pobrany kontekst jest istotny, ani czy odpowiedź
          jest pokryta źródłami --- model może halucynować lub uczciwie zgłosić brak danych, lecz bez gwarancji spójności.
\end{itemize}

\vspace{1em}

% ---------------------------------------------------------------------------
\subsection*{GraphRAG (Hybrid)}

\begin{tcolorbox}[colback=gray!5!white,colframe=gray!60!black,title=Przepływ danych --- pipeline liniowy z dwoma źródłami]
    \centering
    \texttt{Zapytanie} $\rightarrow$ \texttt{Analiza intencji} $\rightarrow$
    \begin{tabular}{c}
        $\swarrow$ \texttt{Graf Neo4j} \\
        $\searrow$ \texttt{Chroma (wektory)}
    \end{tabular}
    $\rightarrow$ \texttt{LLM} $\rightarrow$ \texttt{Odpowiedź}
\end{tcolorbox}

\noindent
Pipeline hybrydowy rozszerza standardowy RAG o warstwę strukturalną.
Najpierw moduł \textit{intent analyzer} klasyfikuje zapytanie (kategoria: terminy, progi, przedmioty maturalne itd.)
i normalizuje encje (nazwy kierunków, tryby studiów). Na tej podstawie uruchamiane są dwa niezależne wyszukiwania:
zapytania Cypher do grafu Neo4j (dane tabelaryczne, relacje między encjami) oraz wyszukiwanie wektorowe po PDF-ach.

\begin{itemize}[leftmargin=*]
    \item \textbf{Zalety:} precyzyjne zapytania strukturalne (filtrowanie numeryczne, mapowanie kierunek$\rightarrow$atrybut),
          łączenie danych z wielu tabel i dokumentów w jednej odpowiedzi, lepsze wyniki przy pytaniach relacyjnych
          (np.\ dyplom zawodowy $\rightarrow$ wskaźnik P2 na danym kierunku).
    \item \textbf{Ograniczenia:} nadal jedno wywołanie generujące --- model otrzymuje \textit{surowy} kontekst z obu źródeł
          bez selekcji; błędne lub nadmiarowe fragmenty mogą zafałszować odpowiedź.
          Graf może gubić logikę tekstową (np.\ operator \textit{albo} w tabeli przedmiotów maturalnych).
    \item \textbf{Brak kontroli:} podobnie jak w RAG wektorowym, nie ma warstwy bezpieczeństwa ani weryfikacji
          pokrycia odpowiedzi w kontekście.
\end{itemize}

\vspace{1em}

% ---------------------------------------------------------------------------
\subsection*{RAG agentowy (LangGraph)}

\begin{tcolorbox}[colback=purple!5!white,colframe=purple!60!black,title=Przepływ danych --- graf stanów z rozgałęzieniami]
    \centering
    \texttt{Zapytanie} $\rightarrow$ \texttt{Guard} $\rightarrow$ \texttt{Intent} $\rightarrow$ \texttt{Hybrid retrieve}
    $\rightarrow$ \texttt{Context filter} $\rightarrow$ \texttt{Answer} $\rightarrow$ \texttt{Fact check}
    $\rightarrow$ \texttt{Odpowiedź / Clarify / Refuse}
\end{tcolorbox}

\noindent
System agentowy nie jest pojedynczym łańcuchem RAG, lecz \textbf{grafem stanów} zbudowanym w LangGraph.
Każdy węzeł to wyspecjalizowany agent (osobne wywołanie LLM lub logika routingu), a krawędzie warunkowe
decydują o dalszym przebiegu przetwarzania. Dzięki temu pipeline może przerwać generowanie, odmówić
odpowiedzi lub poprosić o doprecyzowanie --- zamiast zawsze zwracać tekst wygenerowany na ślepo.

\begin{figure}[H]
    \centering
    \includegraphics[width=0.72\textwidth]{images/agentic-pipeline.png}
    \caption{Diagram aktywności pipeline'u agentowego. Źródło: implementacja w \texttt{agents/graph.py}.}
\end{figure}

\begin{tcolorbox}[colback=blue!5!white,colframe=blue!75!black,title=Opis etapów pipeline'u agentowego]
    \begin{enumerate}[leftmargin=*, nosep]
        \item \textbf{Guarding Agent (strażnik bezpieczeństwa).}
              Pierwszy punkt kontroli. Agent ocenia, czy zapytanie dotyczy spraw uczelnianych (rekrutacja, studia, AGH),
              czy stanowi próbę \textit{jailbreaku}, \textit{prompt injection} albo temat całkowicie niezwiązany z uczelnią.
              \begin{itemize}[nosep]
                  \item \textit{Bezpieczne} $\rightarrow$ przepływ kontynuowany.
                  \item \textit{Jailbreak / injection / off-topic} $\rightarrow$ węzeł \textbf{Refuse} i zakończenie z komunikatem odmowy.
              \end{itemize}

        \item \textbf{Intent Analyzer (analiza intencji).}
              Ten sam moduł normalizacji co w Hybrid RAG: wyciąga kategorię pytania (terminy, progi, matura, kierunek itd.)
              i standaryzuje nazwy kierunków oraz trybów studiów. Wynik kieruje zapytania do grafu Neo4j.

        \item \textbf{Hybrid-retrieve (pobranie hybrydowe).}
              Równoległe pobranie węzłów z grafu wiedzy oraz fragmentów PDF z Chroma (domyślnie do 5 chunków wektorowych
              plus wyniki grafowe). Na tym etapie \textit{nie} następuje jeszcze generacja odpowiedzi --- zbierany jest
              surowy zestaw kandydatów-kontekstów.

        \item \textbf{Context Filtering Agent (filtr kontekstu).}
              LLM ocenia \textbf{każdy} pobrany fragment pod kątem istotności względem pytania i rozpoznanej intencji.
              Fragmenty oznaczone jako \texttt{reject} są odrzucane (np.\ cykle rekrutacyjne 2 i 3, gdy pytanie dotyczy cyklu 1).
              \begin{itemize}[nosep]
                  \item Brak fragmentów po filtracji $\rightarrow$ \textbf{Clarify template} (prośba o doprecyzowanie / informacja o braku danych).
                  \item Pozostały istotne fragmenty $\rightarrow$ przejście do generacji.
              \end{itemize}

        \item \textbf{Answer Generating Agent (generacja odpowiedzi).}
              Model tworzy \textit{szkic odpowiedzi} (\texttt{draft\_answer}) wyłącznie na podstawie przefiltrowanego kontekstu.
              To nadal etap generatywny, lecz oparty na znacznie ciaśniejszym, jakościowo ocenionym zbiorze źródeł.

        \item \textbf{Fact Check Agent (weryfikacja faktów).}
              Osobny agent sprawdza, czy \textbf{każde istotne twierdzenie} w szkicu ma pokrycie w dostarczonym kontekście.
              Nie ocenia stylu ani kompletności bazy --- jedynie wykrywa halucynacje i niespójności liczbowe.
              \begin{itemize}[nosep]
                  \item \textit{Grounded = true} $\rightarrow$ szkic staje się odpowiedzią końcową (\textbf{Send answer}).
                  \item \textit{Grounded = false} $\rightarrow$ powrót do \textbf{Clarify template} z opisem wykrytych problemów
                        (np.\ błędna suma punktów sportowych).
              \end{itemize}
    \end{enumerate}
\end{tcolorbox}

\vspace{0.5em}

\noindent
Kluczowa różnica względem pozostałych architektur: system agentowy stosuje \textbf{wielokrotne, celowe wywołania LLM}
(zwykle 4--6 na jedno pytanie) z jawnymi punktami decyzyjnymi. Zamiast jednego monolitycznego promptu otrzymujemy
łańcuch wyspecjalizowanych ról: strażnik, selektor kontekstu, generator, audytor. Latencja rośnie, lecz rośnie też
możliwość odrzucenia błędnej odpowiedzi zanim trafi do użytkownika.

\vspace{1em}

% ---------------------------------------------------------------------------
\subsection*{Porównanie architektur}

\noindent
\begin{tabularx}{\textwidth}{@{} >{\raggedright\arraybackslash}p{2.4cm} >{\raggedright\arraybackslash}X >{\raggedright\arraybackslash}X >{\raggedright\arraybackslash}X @{}}
    \toprule
    \textbf{Aspekt} & \textbf{Standardowy RAG} & \textbf{GraphRAG (Hybrid)} & \textbf{RAG agentowy} \\
    \midrule
    Źródła danych &
    Tylko PDF (Chroma) &
  Graf Neo4j + PDF &
  Graf Neo4j + PDF \\
    \addlinespace
    Liczba wywołań LLM &
    1 &
  1 &
  4--6 (guard, filter, answer, fact-check) \\
    \addlinespace
    Analiza intencji &
    Nie &
  Tak (reguły + normalizacja) &
  Tak (ten sam moduł) \\
    \addlinespace
    Filtrowanie kontekstu &
    Brak (top-$k$ embeddingów) &
  Brak (cały retrieval trafia do promptu) &
  Tak (LLM ocenia każdy fragment) \\
    \addlinespace
    Weryfikacja odpowiedzi &
    Brak &
  Brak &
  Tak (fact-check przed wysłaniem) \\
    \addlinespace
    Bezpieczeństwo &
    Brak &
  Brak &
  Tak (guard: jailbreak, injection, off-topic) \\
    \addlinespace
    Ścieżki wyjścia &
    Zawsze odpowiedź LLM &
  Zawsze odpowiedź LLM &
  Odpowiedź / odmowa / prośba o doprecyzowanie \\
    \addlinespace
    Latencja &
    Najniższa &
  Średnia &
  Najwyższa \\
    \bottomrule
\end{tabularx}

\vspace{1em}

\begin{tcolorbox}[colback=green!5!white,colframe=green!50!black,title=Wnioski projektowe]
    \begin{itemize}[leftmargin=*]
        \item \textbf{Komplementarność retrieval i kontroli:} Hybrid RAG rozwiązuje problem \textit{dostępu} do właściwych danych
              (struktura grafu + semantyka wektorów), natomiast RAG agentowy dodaje warstwę \textit{jakości i bezpieczeństwa} ---
              selekcję kontekstu, audyt odpowiedzi oraz blokadę zapytań spoza domeny.
        \item \textbf{Koszt za jakość:} Każdy dodatkowy agent to osobne wywołanie modelu. W benchmarku latencja pipeline'u agentowego
              jest wielokrotnie wyższa niż w wariantach liniowych, co jest akceptowalnym kompromisem w scenariuszach,
              gdzie priorytetem jest wiarygodność (rekrutacja, progi punktowe, regulaminy).
        \item \textbf{Fail-safe zamiast fail-open:} Gdy filtr nie znajdzie istotnych fragmentów lub fact-check wykryje halucynację,
              system \textbf{nie zgaduje} --- zwraca szablon doprecyzowania lub komunikat o braku pokrycia w bazie wiedzy.
              To zachowanie widoczne m.in.\ przy pytaniach o punkty sportowe, gdzie agent odrzucił błędny szkic zamiast
              podać niepoprawną maksymalną liczbę punktów.
        \item \textbf{Ochrona przed atakami:} Warstwa Guard skutecznie blokuje zapytania typu jailbreak i prompt injection
              (kategoria \texttt{security} w zestawie testowym), czego liniowe pipeline'y w ogóle nie implementują.
    \end{itemize}
\end{tcolorbox}

\vspace{1em}

\newpage

% Wymaga pakietów z main.tex: tcolorbox, graphicx, booktabs, tabularx, enumitem

\section*{Architektura systemów RAG}

W ramach benchmarku porównano trzy podejścia do odpowiadania na pytania dotyczące rekrutacji i studiów na AGH.
Wszystkie warianty korzystają z tej samej bazy wiedzy (dokumenty PDF w Chroma oraz graf Neo4j), lecz różnią się
sposobem doboru kontekstu, liczbą wywołań modelu językowego oraz mechanizmami kontroli jakości odpowiedzi.

\begin{tcolorbox}[colback=blue!5!white,colframe=blue!75!black,title=Zakres porównania]
    \begin{itemize}[leftmargin=*, nosep]
        \item \textbf{Standardowy RAG} --- wyszukiwanie wyłącznie wektorowe (Chroma) + jedno wywołanie LLM.
        \item \textbf{GraphRAG (Hybrid)} --- równoległy retrieval z grafu wiedzy i wektorów + jedno wywołanie LLM.
        \item \textbf{RAG agentowy} --- orkiestracja wieloetapowa (LangGraph) z warstwą bezpieczeństwa, filtrem kontekstu i weryfikacją faktów.
    \end{itemize}
\end{tcolorbox}

\vspace{1em}

% ---------------------------------------------------------------------------
\subsection*{Standardowy RAG (wektorowy)}

\begin{tcolorbox}[colback=gray!5!white,colframe=gray!60!black,title=Przepływ danych --- pipeline liniowy]
    \centering
    \texttt{Zapytanie} $\rightarrow$ \texttt{Embedding} $\rightarrow$ \texttt{Chroma (top-$k$)} $\rightarrow$ \texttt{LLM} $\rightarrow$ \texttt{Odpowiedź}
\end{tcolorbox}

\noindent
Standardowy RAG implementuje klasyczny schemat \textit{Retrieve--Augment--Generate}.
Zapytanie użytkownika jest embedowane modelem \texttt{nomic-embed-text} i porównywane z fragmentami
PDF-ów zgromadzonymi w bazie Chroma. Trzy najbardziej podobne semantycznie chunki trafiają do jednego
promptu, na podstawie którego model generuje odpowiedź.

\begin{itemize}[leftmargin=*]
    \item \textbf{Zalety:} prostota, niska latencja (jedno wywołanie LLM), dobre wyniki przy pytaniach,
          gdy właściwy fragment tekstu leży blisko zapytania w przestrzeni embeddingów.
    \item \textbf{Ograniczenia:} brak rozumienia struktury danych (tabele, progi, relacje kierunek$\rightarrow$przedmiot),
          podatność na \textit{semantic mismatch} --- trafienie w dokument o podobnym słownictwie, lecz innej treści
          (np.\ punkty ECTS zamiast progów rekrutacyjnych).
    \item \textbf{Brak kontroli:} system nie weryfikuje, czy pobrany kontekst jest istotny, ani czy odpowiedź
          jest pokryta źródłami --- model może halucynować lub uczciwie zgłosić brak danych, lecz bez gwarancji spójności.
\end{itemize}

\vspace{1em}

% ---------------------------------------------------------------------------
\subsection*{GraphRAG (Hybrid)}

\begin{tcolorbox}[colback=gray!5!white,colframe=gray!60!black,title=Przepływ danych --- pipeline liniowy z dwoma źródłami]
    \centering
    \texttt{Zapytanie} $\rightarrow$ \texttt{Analiza intencji} $\rightarrow$
    \begin{tabular}{c}
        $\swarrow$ \texttt{Graf Neo4j} \\
        $\searrow$ \texttt{Chroma (wektory)}
    \end{tabular}
    $\rightarrow$ \texttt{LLM} $\rightarrow$ \texttt{Odpowiedź}
\end{tcolorbox}

\noindent
Pipeline hybrydowy rozszerza standardowy RAG o warstwę strukturalną.
Najpierw moduł \textit{intent analyzer} klasyfikuje zapytanie (kategoria: terminy, progi, przedmioty maturalne itd.)
i normalizuje encje (nazwy kierunków, tryby studiów). Na tej podstawie uruchamiane są dwa niezależne wyszukiwania:
zapytania Cypher do grafu Neo4j (dane tabelaryczne, relacje między encjami) oraz wyszukiwanie wektorowe po PDF-ach.

\begin{itemize}[leftmargin=*]
    \item \textbf{Zalety:} precyzyjne zapytania strukturalne (filtrowanie numeryczne, mapowanie kierunek$\rightarrow$atrybut),
          łączenie danych z wielu tabel i dokumentów w jednej odpowiedzi, lepsze wyniki przy pytaniach relacyjnych
          (np.\ dyplom zawodowy $\rightarrow$ wskaźnik P2 na danym kierunku).
    \item \textbf{Ograniczenia:} nadal jedno wywołanie generujące --- model otrzymuje \textit{surowy} kontekst z obu źródeł
          bez selekcji; błędne lub nadmiarowe fragmenty mogą zafałszować odpowiedź.
          Graf może gubić logikę tekstową (np.\ operator \textit{albo} w tabeli przedmiotów maturalnych).
    \item \textbf{Brak kontroli:} podobnie jak w RAG wektorowym, nie ma warstwy bezpieczeństwa ani weryfikacji
          pokrycia odpowiedzi w kontekście.
\end{itemize}

\vspace{1em}

% ---------------------------------------------------------------------------
\subsection*{RAG agentowy (LangGraph)}

\begin{tcolorbox}[colback=purple!5!white,colframe=purple!60!black,title=Przepływ danych --- graf stanów z rozgałęzieniami]
    \centering
    \texttt{Zapytanie} $\rightarrow$ \texttt{Guard} $\rightarrow$ \texttt{Intent} $\rightarrow$ \texttt{Hybrid retrieve}
    $\rightarrow$ \texttt{Context filter} $\rightarrow$ \texttt{Answer} $\rightarrow$ \texttt{Fact check}
    $\rightarrow$ \texttt{Odpowiedź / Clarify / Refuse}
\end{tcolorbox}

\noindent
System agentowy nie jest pojedynczym łańcuchem RAG, lecz \textbf{grafem stanów} zbudowanym w LangGraph.
Każdy węzeł to wyspecjalizowany agent (osobne wywołanie LLM lub logika routingu), a krawędzie warunkowe
decydują o dalszym przebiegu przetwarzania. Dzięki temu pipeline może przerwać generowanie, odmówić
odpowiedzi lub poprosić o doprecyzowanie --- zamiast zawsze zwracać tekst wygenerowany na ślepo.

\begin{figure}[H]
    \centering
    \includegraphics[width=0.72\textwidth]{images/agentic-pipeline.png}
    \caption{Diagram aktywności pipeline'u agentowego. Źródło: implementacja w \texttt{agents/graph.py}.}
\end{figure}

\begin{tcolorbox}[colback=blue!5!white,colframe=blue!75!black,title=Opis etapów pipeline'u agentowego]
    \begin{enumerate}[leftmargin=*, nosep]
        \item \textbf{Guarding Agent (strażnik bezpieczeństwa).}
              Pierwszy punkt kontroli. Agent ocenia, czy zapytanie dotyczy spraw uczelnianych (rekrutacja, studia, AGH),
              czy stanowi próbę \textit{jailbreaku}, \textit{prompt injection} albo temat całkowicie niezwiązany z uczelnią.
              \begin{itemize}[nosep]
                  \item \textit{Bezpieczne} $\rightarrow$ przepływ kontynuowany.
                  \item \textit{Jailbreak / injection / off-topic} $\rightarrow$ węzeł \textbf{Refuse} i zakończenie z komunikatem odmowy.
              \end{itemize}

        \item \textbf{Intent Analyzer (analiza intencji).}
              Ten sam moduł normalizacji co w Hybrid RAG: wyciąga kategorię pytania (terminy, progi, matura, kierunek itd.)
              i standaryzuje nazwy kierunków oraz trybów studiów. Wynik kieruje zapytania do grafu Neo4j.

        \item \textbf{Hybrid-retrieve (pobranie hybrydowe).}
              Równoległe pobranie węzłów z grafu wiedzy oraz fragmentów PDF z Chroma (domyślnie do 5 chunków wektorowych
              plus wyniki grafowe). Na tym etapie \textit{nie} następuje jeszcze generacja odpowiedzi --- zbierany jest
              surowy zestaw kandydatów-kontekstów.

        \item \textbf{Context Filtering Agent (filtr kontekstu).}
              LLM ocenia \textbf{każdy} pobrany fragment pod kątem istotności względem pytania i rozpoznanej intencji.
              Fragmenty oznaczone jako \texttt{reject} są odrzucane (np.\ cykle rekrutacyjne 2 i 3, gdy pytanie dotyczy cyklu 1).
              \begin{itemize}[nosep]
                  \item Brak fragmentów po filtracji $\rightarrow$ \textbf{Clarify template} (prośba o doprecyzowanie / informacja o braku danych).
                  \item Pozostały istotne fragmenty $\rightarrow$ przejście do generacji.
              \end{itemize}

        \item \textbf{Answer Generating Agent (generacja odpowiedzi).}
              Model tworzy \textit{szkic odpowiedzi} (\texttt{draft\_answer}) wyłącznie na podstawie przefiltrowanego kontekstu.
              To nadal etap generatywny, lecz oparty na znacznie ciaśniejszym, jakościowo ocenionym zbiorze źródeł.

        \item \textbf{Fact Check Agent (weryfikacja faktów).}
              Osobny agent sprawdza, czy \textbf{każde istotne twierdzenie} w szkicu ma pokrycie w dostarczonym kontekście.
              Nie ocenia stylu ani kompletności bazy --- jedynie wykrywa halucynacje i niespójności liczbowe.
              \begin{itemize}[nosep]
                  \item \textit{Grounded = true} $\rightarrow$ szkic staje się odpowiedzią końcową (\textbf{Send answer}).
                  \item \textit{Grounded = false} $\rightarrow$ powrót do \textbf{Clarify template} z opisem wykrytych problemów
                        (np.\ błędna suma punktów sportowych).
              \end{itemize}
    \end{enumerate}
\end{tcolorbox}

\vspace{0.5em}

\noindent
Kluczowa różnica względem pozostałych architektur: system agentowy stosuje \textbf{wielokrotne, celowe wywołania LLM}
(zwykle 4--6 na jedno pytanie) z jawnymi punktami decyzyjnymi. Zamiast jednego monolitycznego promptu otrzymujemy
łańcuch wyspecjalizowanych ról: strażnik, selektor kontekstu, generator, audytor. Latencja rośnie, lecz rośnie też
możliwość odrzucenia błędnej odpowiedzi zanim trafi do użytkownika.

\vspace{1em}

% ---------------------------------------------------------------------------
\subsection*{Porównanie architektur}

\noindent
\begin{tabularx}{\textwidth}{@{} >{\raggedright\arraybackslash}p{2.4cm} >{\raggedright\arraybackslash}X >{\raggedright\arraybackslash}X >{\raggedright\arraybackslash}X @{}}
    \toprule
    \textbf{Aspekt} & \textbf{Standardowy RAG} & \textbf{GraphRAG (Hybrid)} & \textbf{RAG agentowy} \\
    \midrule
    Źródła danych &
    Tylko PDF (Chroma) &
  Graf Neo4j + PDF &
  Graf Neo4j + PDF \\
    \addlinespace
    Liczba wywołań LLM &
    1 &
  1 &
  4--6 (guard, filter, answer, fact-check) \\
    \addlinespace
    Analiza intencji &
    Nie &
  Tak (reguły + normalizacja) &
  Tak (ten sam moduł) \\
    \addlinespace
    Filtrowanie kontekstu &
    Brak (top-$k$ embeddingów) &
  Brak (cały retrieval trafia do promptu) &
  Tak (LLM ocenia każdy fragment) \\
    \addlinespace
    Weryfikacja odpowiedzi &
    Brak &
  Brak &
  Tak (fact-check przed wysłaniem) \\
    \addlinespace
    Bezpieczeństwo &
    Brak &
  Brak &
  Tak (guard: jailbreak, injection, off-topic) \\
    \addlinespace
    Ścieżki wyjścia &
    Zawsze odpowiedź LLM &
  Zawsze odpowiedź LLM &
  Odpowiedź / odmowa / prośba o doprecyzowanie \\
    \addlinespace
    Latencja &
    Najniższa &
  Średnia &
  Najwyższa \\
    \bottomrule
\end{tabularx}

\vspace{1em}

\begin{tcolorbox}[colback=green!5!white,colframe=green!50!black,title=Wnioski projektowe]
    \begin{itemize}[leftmargin=*]
        \item \textbf{Komplementarność retrieval i kontroli:} Hybrid RAG rozwiązuje problem \textit{dostępu} do właściwych danych
              (struktura grafu + semantyka wektorów), natomiast RAG agentowy dodaje warstwę \textit{jakości i bezpieczeństwa} ---
              selekcję kontekstu, audyt odpowiedzi oraz blokadę zapytań spoza domeny.
        \item \textbf{Koszt za jakość:} Każdy dodatkowy agent to osobne wywołanie modelu. W benchmarku latencja pipeline'u agentowego
              jest wielokrotnie wyższa niż w wariantach liniowych, co jest akceptowalnym kompromisem w scenariuszach,
              gdzie priorytetem jest wiarygodność (rekrutacja, progi punktowe, regulaminy).
        \item \textbf{Fail-safe zamiast fail-open:} Gdy filtr nie znajdzie istotnych fragmentów lub fact-check wykryje halucynację,
              system \textbf{nie zgaduje} --- zwraca szablon doprecyzowania lub komunikat o braku pokrycia w bazie wiedzy.
              To zachowanie widoczne m.in.\ przy pytaniach o punkty sportowe, gdzie agent odrzucił błędny szkic zamiast
              podać niepoprawną maksymalną liczbę punktów.
        \item \textbf{Ochrona przed atakami:} Warstwa Guard skutecznie blokuje zapytania typu jailbreak i prompt injection
              (kategoria \texttt{security} w zestawie testowym), czego liniowe pipeline'y w ogóle nie implementują.
    \end{itemize}
\end{tcolorbox}

\vspace{1em}

\newpage
